\section*{Chapitre \ref{ch:g6:how-plants-spread} --- Comment les plantes se répandent}

\begin{solution}{exo:g6:how-plants-spread:1}
La \omterm{prop:g6:how-plants-spread:travel}{colonisation} est l'occupation d'un sol nouveau par une \omterm{def:g5:biodiversity:species}{espèce}. La
\omterm{def:g6:how-plants-spread:dispersal}{dissémination} est le transport des \omterm{prop:g3:flowers-fruits-seeds:transform}{graines}, des \omterm{prop:g3:flowers-fruits-seeds:transform}{fruits} ou des \omterm{def:g4:plants-without-seeds:spore}{spores}
loin de la \omterm{def:g1:plants-around-us:plant}{plante} mère~; les transporteurs sont le vent, les \omterm{def:g1:animals-around-us:animal}{animaux}
(au-dehors et au-dedans), l'eau --- et quelques \omterm{prop:g3:flowers-fruits-seeds:transform}{fruits} qui se
projettent eux-mêmes.
\end{solution}

\begin{solution}{exo:g6:how-plants-spread:2}
Au pied de sa mère, elle germerait à son ombre et parmi ses \omterm{def:g1:plants-around-us:parts}{racines},
en se disputant la \omterm{def:g6:exploring-our-environment:conditions}{lumière} et l'eau avec une rivale cent fois plus
grande. La distance échappe à ce match perdu d'avance.
\end{solution}

\begin{solution}{exo:g6:how-plants-spread:3}
Aigrette du pissenlit~: le vent. \omterm{prop:g3:flowers-fruits-seeds:transform}{Fruit} de la bardane~: les poils des
\omterm{def:g1:animals-around-us:animal}{animaux}, au-dehors. Cerise~: le \omterm{def:g4:journey-of-food:tract}{tube digestif} d'un \omterm{def:g1:animals-around-us:animal}{animal},
au-dedans. Samare d'érable~: le vent (une aile qui tournoie). Noix de
coco~: l'eau.
\end{solution}

\begin{solution}{exo:g6:how-plants-spread:4}
Leur enveloppe coriace leur permet de traverser la \omterm{def:g4:journey-of-food:digestion}{digestion} de
l'oiseau sans dommage~; elles arrivent déposées loin de chez elles,
avec une dose d'engrais.
\end{solution}

\begin{solution}{exo:g6:how-plants-spread:5}
La projection~: leurs gousses sèchent, se tendent comme des ressorts
et claquent, projetant les \omterm{prop:g3:flowers-fruits-seeds:transform}{graines} à plusieurs mètres de la \omterm{def:g1:plants-around-us:plant}{plante}
mère.
\end{solution}

\begin{solution}{exo:g6:how-plants-spread:6}
Le chiendent, qui avance sous terre par des \omterm{def:g1:plants-around-us:parts}{tiges} rampantes~; les
ronces, qui s'arquent et s'enracinent par la pointe (ou le fraisier,
qui saute par ses \omterm{ex:g4:plants-without-seeds:runner}{stolons}).
\end{solution}

\begin{solution}{exo:g6:how-plants-spread:7}
Pissenlit~: arrivé par le vent grâce à ses aigrettes. Coquelicot~:
ses \omterm{prop:g3:flowers-fruits-seeds:transform}{graines} étaient déjà dans le sol, réveillées quand la terre a
été remuée. Mûrier~: déposé par les \omterm{prop:g3:sorting-living-things:groups}{oiseaux} perchés, \omterm{prop:g3:flowers-fruits-seeds:transform}{graines} à
l'intérieur des fientes. Bardane~: apportée accrochée aux poils d'un
chien.
\end{solution}

\begin{solution}{exo:g6:how-plants-spread:8}
(a) Équipement de vol --- des aigrettes~: le vent. (b) Charnu,
coloré, \omterm{prop:g3:flowers-fruits-seeds:transform}{graines} à enveloppe coriace à l'intérieur~: mangé par des
\omterm{def:g1:animals-around-us:animal}{animaux}, les \omterm{prop:g3:flowers-fruits-seeds:transform}{graines} voyagent dans le \omterm{def:g4:journey-of-food:tract}{tube digestif}. (c) Des
crochets qui s'agrippent à la laine~: les poils des \omterm{def:g1:animals-around-us:animal}{animaux} --- et
l'essai sur le pull vient de le confirmer.
\end{solution}

\begin{solution}{exo:g6:how-plants-spread:9}
Les \omterm{def:g1:plants-around-us:parts}{tiges} arquées qui s'enracinent par la pointe font le mètre ---
une occupation dense et sûre du terrain d'à côté. Les baies mangées
par les \omterm{prop:g3:sorting-living-things:groups}{oiseaux} font le kilomètre --- de rares atterrissages
lointains qui fondent de nouveaux buissons.
\end{solution}

\begin{solution}{exo:g6:how-plants-spread:10}
Les \omterm{def:g4:plants-without-seeds:spore}{spores} atteignent les deux talus --- elles flottent partout ---
mais la reproduction d'une \omterm{def:g1:plants-around-us:plant}{plante} à \omterm{def:g4:plants-without-seeds:spore}{spores} exige de l'\omterm{def:g6:exploring-our-environment:conditions}{humidité}, si
bien que seules les \omterm{def:g6:exploring-our-environment:conditions}{conditions} du talus nord acceptent l'arrivée. La
\omterm{def:g6:how-plants-spread:dispersal}{dissémination} livre~; les \omterm{def:g6:exploring-our-environment:conditions}{conditions} décident qui peut rester~: les
habitants correspondent aux \omterm{def:g6:exploring-our-environment:conditions}{conditions}, quelles que soient les
arrivées.
\end{solution}

\begin{solution}{exo:g6:how-plants-spread:11}
Seuls les services à longue portée~: le vent (\omterm{prop:g3:flowers-fruits-seeds:transform}{graines} plumeuses ou
fines comme de la poussière, \omterm{def:g4:plants-without-seeds:spore}{spores}), l'eau de mer (\omterm{prop:g3:flowers-fruits-seeds:transform}{fruits} flottants
comme la noix de coco) et les \omterm{prop:g3:sorting-living-things:groups}{oiseaux} (\omterm{prop:g3:flowers-fruits-seeds:transform}{graines} de \omterm{prop:g3:flowers-fruits-seeds:transform}{fruits} charnus
portées dans le \omterm{def:g4:journey-of-food:tract}{tube digestif}). Prédis chez les pionnières un
équipement plumeux, poussiéreux, flottant ou porté par des baies ---
aucune \omterm{def:g2:from-seed-to-plant:seed}{graine} lourde et sans transporteur.
\end{solution}

\begin{solution}{exo:g6:how-plants-spread:12}
Le sucre achète du transport~: un oiseau payé en chair emporte les
\omterm{prop:g3:flowers-fruits-seeds:transform}{graines} loin de l'ombre de la mère et les dépose avec de l'engrais
--- la distance, ce premier cadeau. L'enveloppe coriace en est
l'autre moitié parce que le messager est un \omterm{def:g4:journey-of-food:tract}{tube digestif}~: le
paiement doit être digestible et les passagères ne doivent pas
l'être. Le sucré recrute le transporteur~; l'enveloppe lui survit.
\end{solution}

\begin{solution}{pb:g6:how-plants-spread:1}
\textbf{1.} Coquelicots~: réveillés de \omterm{prop:g3:flowers-fruits-seeds:transform}{graines} déjà présentes dans
la terre labourée --- le stock en attente des \omterm{prop:g6:life-through-seasons:herbs}{annuelles}. Pissenlits
et chardons~: arrivés par le vent, grâce à leurs \omterm{prop:g3:flowers-fruits-seeds:transform}{graines} plumeuses.

\textbf{2.} Les \omterm{prop:g3:sorting-living-things:groups}{oiseaux} qui mangent des mûres ailleurs se perchent
sur les piquets et déposent les \omterm{prop:g3:flowers-fruits-seeds:transform}{graines} en dessous, engrais compris
--- les plantules cartographient donc les perchoirs, pas le vent.

\textbf{3.} Sa sécheresse~: un champ découvert et ensoleillé n'offre
rien de l'\omterm{def:g6:exploring-our-environment:conditions}{humidité} qu'exigent les stades délicats issus des \omterm{def:g4:plants-without-seeds:spore}{spores}
de fougère. L'arrivée a eu lieu~; c'est l'installation qui a été
refusée.

\textbf{4.} L'avancée végétative --- des \omterm{def:g1:plants-around-us:parts}{tiges} rampantes
souterraines. Elle ne progresse que de proche en proche à partir de
\omterm{def:g1:plants-around-us:plant}{plantes} établies, et doit donc entrer par la bordure~; elle n'a
aucun stade aérien pour se poser au milieu.

\textbf{5.} Dans le sol, sous forme de \omterm{prop:g3:flowers-fruits-seeds:transform}{graines}~: les \omterm{prop:g6:life-through-seasons:herbs}{annuelles}
battues dans la compétition sur pied persistent comme banque de
\omterm{prop:g3:flowers-fruits-seeds:transform}{graines}, en attente du prochain sol nu.

\textbf{6.} Le service au-dedans des \omterm{def:g1:animals-around-us:animal}{animaux} (baies mangées, \omterm{prop:g3:flowers-fruits-seeds:transform}{graines}
déposées depuis les perchoirs --- plus denses sous la ligne de
clôture et les haies) et, pour l'aubépine tout autant, le même
messager à \omterm{def:g1:animals-around-us:covering}{plumes}~; la répartition en forme de perchoirs est
l'indice.

\textbf{7.} Les glands de chêne et les marrons sont lourds et sans
transporteur si personne ne les aide --- ils tombent au pied de leur
mère. Leur messager est le geai (et l'écureuil)~: en enterrant des
glands comme réserves d'hiver et en en oubliant quelques-uns, ils
plantent des chênes loin de la mère.

\textbf{8.} La végétation de chaque année ombrage le sol en dessous,
le refroidit et l'obscurcit~: la réussite même des pionnières
amatrices de soleil supprime leurs \omterm{def:g6:exploring-our-environment:conditions}{conditions}, et les \omterm{def:g6:exploring-our-environment:conditions}{conditions}
changées favorisent ensuite les amateurs d'ombre --- les colons
modifient le site, et le site modifié choisit de nouveaux colons.

\textbf{9.} Vague 1, les \omterm{prop:g3:flowers-fruits-seeds:transform}{graines} en attente dans le sol
(instantanée --- elles étaient déjà là). Vague 2, les voyageurs du
vent (des semaines~: l'air livre en permanence). Vague 3, les
arbustes portés par les \omterm{def:g1:animals-around-us:animal}{animaux} (des années~: perchoir après
perchoir). Vague 4, les arbres à \omterm{prop:g3:flowers-fruits-seeds:transform}{graines} lourdes (en dernier~: ils
attendent les oublis enterrés des geais et des écureuils).

\textbf{10.} Les fissures n'offrent pas de place~: les \omterm{def:g1:plants-around-us:parts}{racines} d'un
arbuste demandent un volume de terre qu'un joint de pavage n'a pas~;
et les vagues suivantes de la succession ont besoin de leurs
messagers et de leurs \omterm{def:g6:exploring-our-environment:conditions}{conditions} --- les dépôts des \omterm{prop:g3:sorting-living-things:groups}{oiseaux} perchés
manquent surtout les fissures, et le joint sec et brûlant refuse ce
qui arrive quand même. Les voyageurs du vent de la première année
sont le seul service que le site peut accepter.

\textbf{11.} La tendance va de la terre nue aux herbes, aux
arbustes, aux arbres pionniers --- chaque vague plus haute et plus
ombreuse~: prolongée, elle finit en bois, comme les saules et les
bouleaux le montrent déjà. L'exception~: un pré fauché ou pâturé, où
la main humaine supprime la vague des arbustes chaque année ---
preuve que ce n'est que \emph{laissé tranquille} que le champ marche
vers la forêt.

\textbf{12.} Par exemple~: «~La \omterm{def:g6:how-plants-spread:dispersal}{dissémination} a livré des colons au
champ nu --- venus du sol, du vent, des \omterm{prop:g3:sorting-living-things:groups}{oiseaux} --- et à chaque
étape les \omterm{def:g6:exploring-our-environment:conditions}{conditions} du champ ont décidé quelles arrivées pouvaient
rester. La \omterm{prop:g6:how-plants-spread:travel}{colonisation} s'est ainsi déroulée par vagues, chaque
vague changeant les \omterm{def:g6:exploring-our-environment:conditions}{conditions} et choisissant la suivante.~»
\end{solution}
